\documentclass[10pt,a4paper]{amsart}
\usepackage[margin=19mm]{geometry}
\usepackage{amsmath,amssymb,mathtools}
\usepackage[scr=rsfs,cal=euler]{mathalfa}
\usepackage{xfrac}
\usepackage[unicode,hidelinks,bookmarks=false]{hyperref}
\newcommand{\ie}{i.e.\ }
\newcommand{\cf}{cf.\ }
\newcommand{\resp}{respectively\ }
\newcommand{\Subsection}{\subsection}
\providecommand{\nobreakdash}{\nobreak\hbox{-}}
\setlength{\emergencystretch}{2em}
\multlinegap0pt
\usepackage{amssymb}
\usepackage{mathtools}
\usepackage[shortlabels]{enumitem}
\newtheorem{lemma}{Lemma}
\newtheorem{theorem}{Theorem}
\newcommand{\Tr}{\operatorname{Tr}}
\newcommand{\norm}[1]{\left\lVert #1\right\rVert}
\title{A negative exponent range for Audenaert's complementary McCarthy trace inequality}
\author{Xing Li, Bin Zhou}
\date{Source: arXiv:2608.20937v1 (CC BY 4.0)}
\begin{document}
\maketitle

\noindent\emph{Source and licence.} A negative exponent range for Audenaert's complementary McCarthy trace inequality. Version arXiv:2608.20937v1 (CC BY 4.0), distributed by arXiv under the Creative Commons Attribution 4.0 International licence, http://creativecommons.org/licenses/by/4.0/, as recorded in the arXiv licence metadata for this exact version and shown on the abstract page https://arxiv.org/abs/2608.20937v1. Full licence notice: \texttt{licenses/arxiv-2608.20937-CC-BY-4.0.txt}.\par\smallskip

\noindent\textit{Editorial scope.} Counted unit: the article's theorem thm:parallel-asym (source0.tex lines 303--311). The article's complete original proof of it is delivered with it (source0.tex lines 313--374, one \texttt{proof} environment of 62 lines). No mathematics is written, repaired or re-derived: the statement and the proof are the article's own lines, in the article's own notation, and no retained line is substituted or re-flowed.  Retained uncounted as the source-local context the proof needs: lines 144--146; lines 155--168; lines 186--192; lines 212--231; lines 218--222.  Nothing was omitted from any retained span, so the package carries no omission record.  No retained line contains a citation, so the wrapper carries no bibliography and no bibliography entry.  No retained line contains a figure environment, an image-inclusion command or drawing code, so no image is delivered and the package declares no asset dependency.  Editorial formatting only, in addition: an AMS \texttt{amsart} preamble that restates the article's own declarations and macros used by the retained lines, namely the environments \texttt{lemma}, \texttt{theorem} on separate counters, together with the standard packages those lines need.  The retained environments print in this document as Lemma 1, Theorem 1 in order of appearance, whereas the article prints the same items with the numbering of its own full document.  The complete native source of the article as distributed in the arXiv e-print for this exact version is bundled unchanged as \texttt{sources/208213/source0.tex}.
\begin{equation}\label{eq:weighted-mean-chain}
    X!_tY\le X\#_tY\le (1-t)X+tY.
\end{equation}

\begin{lemma}[Geometric means and products]\label{lem:ando-hiai}
Let \(X,Y>0\) and \(0<t<1\).  Then
\begin{equation}\label{eq:AH-weighted}
    \lambda(X\#_tY)\prec_{\log}s(X^{1-t}Y^t).
\end{equation}
Consequently, for every \(r>0\),
\begin{equation}\label{eq:geomean-trace-bound-weighted}
    \Tr(X\#_tY)^r\le \norm{X^{1-t}Y^t}_r^r.
\end{equation}
In particular, when \(t=1/2\),
\begin{equation}\label{eq:geomean-trace-bound}
    \Tr(X\#Y)^r\le \norm{X^{1/2}Y^{1/2}}_r^r.
\end{equation}
\end{lemma}

\begin{lemma}[Trace monotonicity and equality]\label{lem:trace-monotone}
If \(0\le C\le D\), then for every \(r>0\),
\[
    \Tr C^r\le \Tr D^r.
\]
Moreover, if equality holds, then \(C=D\).
\end{lemma}

\begin{lemma}[Finite-dimensional Schatten H\"older inequality]\label{lem:holder}
Let \(p,q,r>0\) and \(1/r=1/p+1/q\).  For all compatible matrices \(U,V\),
\[
    \norm{UV}_r\le \norm{U}_p\norm{V}_q.
\]
In particular, for \(X,Y>0\), \(r>0\), and \(0<t<1\),
\begin{equation}\label{eq:holder-weighted-special}
    \norm{X^{1-t}Y^t}_r^r
    \le
    (\Tr X^r)^{1-t}(\Tr Y^r)^t.
\end{equation}
At \(t=1/2\),
\begin{equation}\label{eq:holder-special}
    \norm{X^{1/2}Y^{1/2}}_r^r
    \le
    (\Tr X^r)^{1/2}(\Tr Y^r)^{1/2}
    \le
    \frac{\Tr X^r+\Tr Y^r}{2}.
\end{equation}
\end{lemma}

\begin{equation}
    \norm{X^{1-t}Y^t}_r^r
    \le
    (\Tr X^r)^{1-t}(\Tr Y^r)^t.
\end{equation}

\begin{theorem}[A one-parameter parallel sum inequality]\label{thm:parallel-asym}
Let \(X,Y>0\), \(r>0\), and \(0<t<1\).  Then
\begin{equation}\label{eq:parallel-asym}
    \Tr(X:Y)^r+
    \bigl(m_t^{-1}-m_t^r\bigr)\norm{X^{1-t}Y^t}_r^r
    \le
    \Tr X^r+\Tr Y^r .
\end{equation}
\end{theorem}

\begin{proof}
Put
\[
    H=X:Y,\qquad
    G_t=\norm{X^{1-t}Y^t}_r^r,\qquad
    S=\Tr X^r+\Tr Y^r.
\]
Since
\[
    (1-t)X^{-1}+tY^{-1}
    \le
    m_t(X^{-1}+Y^{-1}),
\]
inverting gives
\[
    X:Y\le m_t(X!_tY).
\]
By \eqref{eq:weighted-mean-chain},
\[
    X:Y\le m_t(X\#_tY).
\]
Using Lemmas \ref{lem:trace-monotone} and \ref{lem:ando-hiai}, we get
\begin{equation}\label{eq:asym-H-estimate}
    \Tr H^r
    \le
    m_t^r\Tr(X\#_tY)^r
    \le
    m_t^r G_t.
\end{equation}
On the other hand, Lemma \ref{lem:holder}, in the form \eqref{eq:holder-weighted-special}, gives
\begin{equation}\label{eq:asym-holder}
    G_t
    =
    \norm{X^{1-t}Y^t}_r^r
    \le
    (\Tr X^r)^{1-t}(\Tr Y^r)^t.
\end{equation}
The scalar weighted arithmetic-geometric mean inequality yields
\[
    (\Tr X^r)^{1-t}(\Tr Y^r)^t
    \le
    (1-t)\Tr X^r+t\Tr Y^r
    \le
    m_tS.
\]
Hence
\begin{equation}\label{eq:asym-G-estimate}
    m_t^{-1}G_t\le S.
\end{equation}
Combining \eqref{eq:asym-H-estimate} and \eqref{eq:asym-G-estimate}, we obtain
\[
    \Tr H^r+
    \bigl(m_t^{-1}-m_t^r\bigr)G_t
    \le
    m_t^rG_t+
    \bigl(m_t^{-1}-m_t^r\bigr)G_t
    =
    m_t^{-1}G_t
    \le S.
\]
This is \eqref{eq:parallel-asym}.
\end{proof}

\end{document}
